\documentclass[11pt]{article}

\usepackage[margin=1in]{geometry}
\usepackage[T1]{fontenc}
\usepackage{booktabs}
\usepackage{xcolor}
\usepackage{hyperref}

\hypersetup{
  colorlinks=true,
  linkcolor=blue,
  citecolor=blue,
  urlcolor=blue
}

\newcommand{\reviewcomment}[1]{\par\noindent\textbf{Reviewer comment.} \emph{#1}\par}
\newcommand{\response}[1]{\par\noindent\textbf{Response.} #1\par}
\newcommand{\manuscriptchange}[1]{\par\smallskip\noindent\textcolor{blue}{\textbf{Added or corrected manuscript text:} #1}\par\smallskip}

\title{Response to Review and Camera-Ready Changes\\
\large Unified Coarse-to-Fine Landmark Localization for Multi-Domain Ultrasound Biometry}
\author{Hamza Rasaee \and Hassan Rivaz}
\date{}

\begin{document}
\maketitle

\noindent We thank the reviewers and challenge organizers for their constructive comments. We preserved the reviewed manuscript's wording and structure and made only targeted camera-ready changes. In the revised manuscript, blue text marks the exact inserted or corrected fragments. The reported system is the frozen final-test Docker submission: one trained checkpoint with a shared public DINOv3 ViT-B encoder and internal task-specific modules. It uses neither independently routed checkpoints nor checkpoint ensembling, model averaging, retrieval-based prediction replacement, or manual correction.

\section*{Response to Reviewer NFgs: Publication Quality and Camera-Ready Requirements}

\reviewcomment{Authors must submit the original LaTeX source files, final camera-ready manuscript, and any applicable supplementary materials to the specified email address using the required \texttt{MWM\_\#XX} subject format.}
\response{We prepared the LLNCS source, bibliography, figure assets, and compiled manuscript PDF. The final package will be emailed to \href{mailto:bai_jieyun@126.com}{bai\_jieyun@126.com} with the OpenReview paper number in the required \texttt{MWM\_\#XX} subject line. No separate supplementary document is required for this paper.}

\reviewcomment{The licence-to-publish form must be completed and signed by the corresponding author.}
\response{The corresponding author will download the form from \url{https://ideal-conf.com/downloads}, complete and sign \path{SNCS_ProceedingsPaper_LTP_ST_SN_Switzerland.docx}, and include it in the final package.}

\reviewcomment{Where applicable, authors should appropriately cite publications associated with the event.}
\response{The paper cites Bai et al., \emph{Foundation Model Challenge for Ultrasound Biometry}, Zenodo, MICCAI 2026, DOI \href{https://doi.org/10.5281/zenodo.19736827}{10.5281/zenodo.19736827}. It also cites the organizer-documented JNU-IFM, PSFHS, IUGC, and FUGC dataset or benchmark publications to identify the relevant data lineage. All in-text citation keys were checked against the bibliography.}

\par\smallskip\noindent\textbf{Additional camera-ready checks.}\par
The manuscript retains the supplied LLNCS class and \texttt{splncs04} bibliography style without geometry or font-size overrides. It is 10 pages: manuscript content and acknowledgements end on page 8, and references occupy pages 9--10. The authors' acknowledgement is retained exactly as supplied.

\manuscriptchange{Funded by B2-2004995 (Fonds de recherche du Qu\'ebec – Nature et technologies B2X-363874) and Natural Sciences and Engineering Research Council of Canada (NSERC) Discovery Grant. OpenAI's ChatGPT was used to assist with English language editing in the Introduction, Methods, and Discussion sections. The authors reviewed, verified, and revised all generated text and are responsible for the manuscript content.}

\section*{Response to Reviewer CKh9}

\reviewcomment{The gap between released-annotation MRE and public validation should be analyzed, including possible domain shift, annotation noise, image quality, and practical implications.}
\response{We agree. First, we corrected the released evaluation using the organizer-confirmed vertical cardiac landmark ordering: average MRE changes from 5.64 to 5.70 pixels, primarily because A4C changes from 7.30 to 7.83 pixels. Second, the revised Discussion now states that released MRE and the challenge score are different quantities. Hidden labels are unavailable, so a per-task hidden-set decomposition cannot be calculated; we therefore identify possible causes without presenting them as measured effects. The frozen checkpoint improved from 24.77 on public validation to 23.29 on final-test Docker evaluation.}

\manuscriptchange{Released MRE and challenge scores are not directly comparable: 5.70 is an unnormalized landmark distance on released annotations, whereas public and final scores aggregate normalized MRE and parameter MAE on hidden images. Hidden labels prevent a per-task decomposition; plausible contributors include acquisition-domain, image-quality, protocol, and annotation shifts. The improvement from 24.77 public validation to 23.29 final Docker evaluation provides direct evidence of frozen-checkpoint generalization.}

\reviewcomment{Content-aware anatomical retrieval requires cost, failure-mode, gating, and ablation analysis.}
\response{The organizers' final clarification excludes routing to independently trained task- or dataset-specific checkpoints from the single-model setting. Retrieval-based prediction replacement is therefore not part of the camera-ready method or final Docker result. The manuscript reports only the compliant one-checkpoint model and does not claim the earlier retrieval-assisted public result.}

\reviewcomment{A structured ablation should isolate offset refinement, retrieval, and heatmap resolution.}
\response{Retrieval is excluded for the compliance reason above. We retain the compact single-checkpoint comparison in Table~1 and correct all released-set values under the organizer-confirmed ordering. The matched resolution comparison shows 5.70 versus 5.59 released MRE for 128- and 256-resolution heatmaps, while the 128-resolution model has the better public score (24.77 versus 24.87). Because no matched one-factor run exists for every simultaneously introduced component, we explicitly avoid claiming unsupported isolated gains.}

\begin{center}
\small
\textbf{Table 1.} Comparison of representative single-checkpoint model variants. ``Rel.'' denotes released annotated data, ``Val.'' denotes the held-out checkpoint-selection split, ``Public'' denotes the challenge public validation score, and ``Final'' denotes the final-test Docker evaluation score.

\vspace{0.4em}
\setlength{\tabcolsep}{3pt}
\begin{tabular}{lcccc}
\toprule
Variant & Rel. MRE$\downarrow$ & Val. MRE$\downarrow$ & Public$\downarrow$ & Final$\downarrow$ \\
\midrule
Early ViT-L task-FPN baseline & \textcolor{blue}{11.02} & 25.80 & 29.88 & -- \\
High-resolution DINOv3 offset model & \textcolor{blue}{5.70} & 15.62 & \textbf{24.77} & \textbf{23.29} \\
256-resolution heatmap variant & \textcolor{blue}{5.59} & 15.96 & 24.87 & -- \\
\bottomrule
\end{tabular}
\end{center}

\reviewcomment{The paper should compare the single-model method against multi-model approaches.}
\response{The challenge clarification states that independently trained task- or dataset-specific checkpoints and checkpoint ensembles do not satisfy the single-model setting. We therefore do not present a non-compliant multi-model pipeline as a directly comparable submission. Instead, the manuscript now quantifies allowed internal specialization: all tasks share one encoder representation and exported checkpoint, while task identity activates internal FPN and decoder modules.}

\reviewcomment{The offset refinement implementation should specify the predicted correction, decoding location, and supervision.}
\response{We added the requested implementation detail immediately after the existing soft-argmax and offset equations. Generic heads predict a bounded two-dimensional correction after differentiable soft-argmax, and direct coordinate loss supervises the refined output. Figure~1 now also distinguishes generic and A4C offset refinement from the geometry-specific FUGC, HC, and IVC refiners.}

\manuscriptchange{For generic offset heads, FPN features sampled at each soft-argmax coordinate are combined with coordinate, global-image, and landmark-index embeddings. Two graph-refinement blocks predict a bounded two-dimensional correction through a \texttt{tanh} output, with maximum normalized displacement 0.03; direct coordinate loss supervises the refined output.}

\section*{Response to Reviewer 3zvF}

\reviewcomment{Analyze which tasks, preprocessing steps, or split differences drive the released-to-public gap.}
\response{The revised manuscript reports all nine released-set task errors, held-out checkpoint-selection MRE and measurement MAE, public validation, and final-test Docker evaluation. Per-task hidden errors cannot be computed without hidden labels. The Discussion therefore distinguishes metrics and identifies acquisition-domain, image-quality, protocol, and annotation shifts as hypotheses rather than measured causes.}

\reviewcomment{The retrieval prior should be evaluated on unseen domains and its overfitting risk addressed.}
\response{Retrieval refinement is not part of the compliant camera-ready method or final Docker submission. All retrieval-assisted claims and results are excluded.}

\reviewcomment{The extent and cost of parameter sharing across nine specialized FPNs, adapters, and heads should be reported.}
\response{We measured the instantiated model. It has 144.4 million trainable parameters: 85.6 million (59.3\%) in the shared DINOv3 encoder and 58.7 million (40.7\%) across all task-specific FPNs, adapters, and heads. All belong to one state dictionary and checkpoint.}

\manuscriptchange{The model has 144.4 million trainable parameters: the shared encoder contains 85.6 million (59.3\%), while all task-specific FPNs, adapters, and heads together contain 58.7 million (40.7\%).}

\reviewcomment{Clarify component contributions, including task-aware losses and 128 versus 256 heatmaps, on the same split.}
\response{The Method already defines the landmark, measurement, and geometry objectives. Table~1 gives the matched 128-versus-256 resolution comparison. We added a limitation to the ablation interpretation rather than assigning unsupported gains where no matched single-factor experiment exists.}

\manuscriptchange{Because no matched one-factor run is available for every component, we do not assign isolated numerical gains to simultaneously introduced changes.}

\reviewcomment{Clinical measurement MAE and failures, especially IVC and cardiac views, need clearer discussion.}
\response{We added the held-out checkpoint-selection measurement MAE of 10.90 pixels, retained the released AOP angle MAE of 3.51 degrees, and report all nine released task MREs. The revised Discussion explicitly identifies IVC, A4C, and HC as the principal released-set failure cases.}

\manuscriptchange{On the held-out checkpoint-selection split, task-averaged MRE and measurement MAE were 15.62 and 10.90 pixels, respectively.}

\manuscriptchange{IVC remains the highest-error released subset and has only 38 samples, while A4C and HC require consistent multi-point geometry. These cases identify limited-sample vessel localization and dense or measurement-sensitive structures as the main remaining failure modes.}

\section*{Additional Reproducibility Corrections}

\response{The final audit also corrected two implementation statements. The selected checkpoint's log confirms that all 6768 released rows were present before splitting, including the 25 fetal-femur images listed later by the organizers as orientation anomalies. The selected 128-resolution model was warm-started from an earlier unified 64-resolution checkpoint trained only on official annotations, after which all 144.4 million parameters remained trainable. These corrections do not change the frozen public or final Docker scores.}

\manuscriptchange{The selected checkpoint was trained on all 6768 released rows available during development, including 25 fetal-femur images later listed by the organizers as orientation anomalies.}

\manuscriptchange{The selected model was warm-started from an earlier unified 64-resolution checkpoint trained only on official annotations; compatible parameters were transferred before all 144.4 million parameters were jointly optimized.}

\section*{Summary}

\begin{itemize}
  \item Preserved the reviewed manuscript and highlighted only exact additions or corrections.
  \item Reported one compliant checkpoint with 24.77 public and 23.29 final-test scores.
  \item Corrected organizer-confirmed released metrics to 5.70 average MRE and 7.83 A4C MRE.
  \item Added metric-gap, offset, parameter-sharing, failure-mode, and reproducibility details.
  \item Removed retrieval-assisted claims from the compliant camera-ready method.
  \item Verified the LLNCS 8+2 page allocation, citations, references, and supplied acknowledgement.
\end{itemize}

\end{document}
